% Part: lambda-calculus
% Chapter: lambda-definability
% Section: primitive-recursive-functions

\documentclass[../../../include/open-logic-section]{subfiles}

\begin{document}

\olfileid{lam}{ldf}{prf}
\olsection{आदिम पुनरावर्ती फलने \usetoken{S}{lambda definable} असतात}

मूलभूत फलने $\Zero$, $\Succ$ आणि $\Proj{n}{i}$
यांपासून संयोजन व आदिम पुनरावर्तन वापरून परिभाषित
करता येणारी फलने म्हणजे आदिम पुनरावर्ती फलने, हे लक्षात घ्या.

\begin{lem}
  \ollabel{lem:basic}
  मूलभूत आदिम पुनरावर्ती फलने $\Zero$, $\Succ$ आणि
  प्रक्षेपणे~$\Proj{n}{i}$ ही !!{lambda definable} आहेत.
\end{lem}

\begin{proof}
पुढील पदे त्यांना !!{lambda define} करतात:
\begin{align*}
  \fn{Zero} & \ident \lambd[a][\lambd[fx][x]]\\
  \fn{Succ} & \ident \lambd[a][\lambd[fx][f (a f x)]] \\
  \fn{Proj}^n_i & \ident \lambd[x_0\dots x_{n-1}][x_i]
\end{align*}
\end{proof}

\begin{lem}
  \ollabel{lem:comp} सर्वत्र परिभाषित $k$-स्थानी फलन $f$ आणि $n$-स्थानी
  सर्वत्र परिभाषित फलने $g_0, \dots, g_{k-1}$ ही अनुक्रमे $F$, $G_0$,
  \dots, $G_{k-1}$ या पदांनी !!{lambda definable} आहेत,
  आणि $h$ हे त्यांच्यापासून संयोजनाने परिभाषित आहे असे समजा.
  मग $h$ हे !!{lambda definable} आहे.
\end{lem}

\begin{proof}
  पुढील पद $h$ ला !!{lambda define} करते:
  \[
  H \ident \lambd[x_0 \dots x_{n-1}][F\, (G_0 x_0 \dots
    x_{n-1}) \dots (G_{k-1} x_0 \dots x_{n-1})]
  \]
  याची पडताळणी सरावासाठी सोडली आहे.
\end{proof}

\begin{prob}
  $H\num{n_0}\dots\num{n_{n-1}} \red \num{h(n_0, \dots,
    n_{n-1})}$ हे दाखवून \olref[lam][ldf][prf]{lem:comp}
  चा पुरावा पूर्ण करा.
\end{prob}

\olref{lem:comp} मध्ये $f$ आणि $g_0$,
\dots,~$g_{k-1}$ ही आदिम पुनरावर्ती असणे आवश्यक
नव्हते, हे लक्षात घ्या; ती सर्वत्र परिभाषित आणि
!!{lambda definable} असणे पुरेसे आहे.

\begin{lem}\ollabel{lem:prim}
  $f$ हे $n$-स्थानी आणि~$g$ हे $n+2$-स्थानी फलन
  असून ती सर्वत्र परिभाषित आणि अनुक्रमे $F$ आणि~$G$ या पदांनी
  !!{lambda definable} आहेत असे समजा. $h$ हे
  $f$ आणि $g$ यांच्यापासून आदिम पुनरावर्तनाने परिभाषित असेल,
  तर $h$ देखील !!{lambda definable} आहे.
\end{lem}

\begin{proof}
  $h$ ची व्याख्या अशी आहे, हे लक्षात घ्या:
  \begin{align*}
    h(x_1, \dots, x_n, 0) &= f(x_1, \dots, x_n)\\
    h(x_1, \dots, x_n, y+1) & = g(x_1, \dots, x_n, y, h(x_1, \dots, x_n, y)).
  \end{align*}
  सहज अर्थाने, आदिम पुनरावर्ती व्याख्येत $f(x_1,
  \dots, x_n)$ या सुरुवातीच्या मूल्यापासून $h$ चे
  मूल्य $y$ वेळा अद्ययावत होते. हे चर्च अंकांच्या
  व्याख्येसारखे आहे: तेही पुनरावर्तक आहेत.

  सोपेपणासाठी आपण एका अतिरिक्त फलसाधक~$x$ साठी
  व्याख्या व पुरावा देऊ. $h$ ला !!{lambda define}
  करणारे पद असे आहे:
  \begin{align*}
    H \ident & \lambd[x][\lambd[y][\fn{Snd} (y
        D \tuple{\num{0}, F x})]]
    \intertext{जेथे }
    D \ident & \lambd[p][\tuple{\fn{Succ} (\fn{Fst}\, p), (G x
          (\fn{Fst}\, p) (\fn{Snd}\, p))}]
  \end{align*}
  पुनरावर्तनात आपण ठेवलेली अवस्था ही जोडी आहे:
  पहिला घटक सध्याचे $y$ आणि दुसरा त्या स्थितीसाठीचे
  $h$ चे मूल्य. प्रत्येक पायरीत $y$ आणि $h$
  यांच्या नव्या मूल्यांची जोडी तयार होते. पुनरावर्तन
  संपल्यावर तिचा दुसरा घटक परत करतो; तेच $h$ चे
  अंतिम मूल्य असते. एकाहून अधिक अतिरिक्त फलसाधक
  असतील, तर ती सर्व निविष्टे सुरुवातीच्या फलनाला आणि
  प्रत्येक पायरीच्या फलनाला त्याच क्रमाने दिली जातात;
  प्रत्येक बाह्य प्राचलासाठी स्वतंत्र अमूर्तीकरण घ्यायचे.
  पुनरावर्तनाच्या जोडीतील मोजणी आणि पूर्वमूल्याचे घटक
  तेच राहतात, म्हणून खालील विगमन तसाच लागू होतो.
  हे निरूपण आदिम पुनरावर्तनाचेच
  आहे, हे आता सिद्ध करू.

  कोणत्याही $n$ आणि $m$ साठी $H\,\num{n}\,\num{m} \red
  \num{h(n,m)}$ दाखवायचे आहे. त्यासाठी आधी $D_n \ident
  \Subst{D}{\num{n}}{x}$ असेल, तर $D_n^m \tuple{\num{0}, F\, \num{n}} \red
  \tuple{\num{m}, \num{h(n, m)}}$ हे दाखवू. $m$ वरील
  विगमनाने पुढे जाऊ.

  $m=0$ असेल, तर $D_n^0 \tuple{\num{0}, F\, \num{n}} \red
  \tuple{\num{0}, \num{h(n, 0)}}$ दाखवायचे आहे.
  पण $D_n^0 \tuple{\num{0}, F\,
    \num{n}}$ म्हणजेच $\tuple{\num{0}, F\, \num{n}}$.
  $F$ हे $f$ ला !!{lambda define}s; त्यामुळे याचे
  $\tuple{\num{0}, \num{f(n)}}$ मध्ये न्यूनीकरण होते;
  आणि $f(n) = h(n, 0)$ असल्याने हेच
  $\tuple{\num{0}, \num{h(n,0)}}$ आहे.

  आता $D_n^m \tuple{\num{0}, F\, \num{n}} \red
  \tuple{\num{m}, \num{h(n, m)}}$ असे समजा.
  $D_n^{m+1} \tuple{\num{0}, F\, \num{n}} \red
  \tuple{\num{m+1}, \num{h(n, m+1)}}$ हे दाखवू:
  \begin{align*}
    D_n^{m+1} \tuple{\num{0}, F\, \num{n}}
    & \ident D_n(D_n^{m} \tuple{\num{0}, F\, \num{n}})\\
    & \red D_n\,\tuple{\num{m}, \num{h(n, m)}} \text{\qquad (विगमन गृहीतकाने)}\\
    & \ident (\lambd[p][
      \tuple{
        \fn{Succ} (\fn{Fst}\, p), (G \, \num{n} (\fn{Fst}\, p) (\fn{Snd}\, p))
    }]) \tuple{\num{m}, \num{h(n,m)}}\\
    & \redone
    \tuple{\fn{Succ} (\fn{Fst}\, \tuple{\num{m}, \num{h(n,m)}}),\\
      & \qquad (G \, \num{n} (\fn{Fst}\, \tuple{\num{m}, \num{h(n,m)}}) (\fn{Snd}\, \tuple{\num{m}, \num{h(n,m)}}))}\\
    & \red
    \tuple{\fn{Succ}\, \num{m}, (G \, \num{n} \,\num{m} \, \num{h(n,m)})}\\
    & \red \tuple{\num{m+1}, \num{g(n, m, h(n, m))}}
  \end{align*}
  $g(n, m, h(n, m)) = h(n, m+1)$ असल्याने हे सिद्ध झाले.

  शेवटी पुढील न्यूनीकरण पाहा:
  \begin{align*}
    H\,\num n\,\num m &\ident \lambd[x][\lambd[y][\fn{Snd} (y
        (\lambd[p].\tuple{\fn{Succ} (\fn{Fst}\, p), (G\, x\,
          (\fn{Fst}\, p)\, (\fn{Snd}\, p))}) \tuple{\num{0}, F x})]]\\
    & \qquad\,\num n\, \num m\\
    & \red \fn{Snd} (\num m \,
    \underbrace{(\lambd[p].\tuple{\fn{Succ} (\fn{Fst}\, p), (G \,\num n\,
      (\fn{Fst}\, p) (\fn{Snd}\, p))})}_{D_n} \tuple{\num{0}, F \num n})\\
    & \ident \fn{Snd} (\num m \, D_n\, \tuple{\num{0}, F \num n})\\
    & \red \fn{Snd} \,(D_n^m  \tuple{\num{0}, F \num n})\\
    & \red \fn{Snd} \,\tuple{\num{m}, \num{h(n, m)}} \\
    & \red \num{h(n,m)}.
  \end{align*}
\end{proof}

\begin{prop}
  प्रत्येक आदिम पुनरावर्ती फलन !!{lambda definable} आहे.
\end{prop}

\begin{proof}
  \olref{lem:basic} नुसार सर्व मूलभूत फलने
  !!{lambda definable} आहेत. तसेच \olref{lem:comp}
  आणि \olref{lem:prim} नुसार सर्वत्र परिभाषित !!{lambda definable}
  फलनांचा वर्ग संयोजन व आदिम पुनरावर्तन यांच्या
  अंतर्गत बंद आहे.
\end{proof}

\end{document}
